\documentclass[11pt,a4paper]{article}
\usepackage[margin=2.5cm]{geometry}
\usepackage{amsmath,amssymb}
\usepackage{booktabs}
\usepackage{listings}
\usepackage[colorlinks=true,linkcolor=blue,citecolor=blue,urlcolor=blue]{hyperref}

\lstset{basicstyle=\ttfamily\small,columns=fullflexible,frame=single,
        language=Fortran,keepspaces=true}

\newcommand{\kb}{\bar k}
\newcommand{\Nc}{\mathcal N}
\newcommand{\E}{\mathcal E}

\title{Rapidity distributions and the initial--initial soft dipole:\\
       partial fractioning of the eikonal in MadNkLO}
\author{MadNkLO internal note}
\date{15 November 2026}

\begin{document}
\maketitle

\begin{abstract}
In the local analytic sector subtraction with initial-state radiation, as
implemented in MadNkLO, the soft dipole between the two incoming partons was
split equally between the two initial--initial momentum mappings (rescaling of
either beam). In the collinear limit along one beam, half of the
collinear-singular part of the soft counterterm was then evaluated on a Born
configuration that is longitudinally boosted, by $\ln z$, relative to the
correct one. Inclusive cross sections and boost-invariant observables were
unaffected, but each bin of a rapidity distribution received a
non-integrable $1/s_{1i}$ term that cancels only between bins, which showed up
as very large Monte Carlo errors. We replace the equal split by a
partial-fraction decomposition of the eikonal, each fraction carrying the
collinear singularity of one beam and being evaluated on the mapping that
rescales that beam. The soft and collinear limits and the integrated
counterterms are unchanged. We report the validation on $pp\to e^+e^-$,
$pp\to Zj$ and $pp\to t\bar t$.
\end{abstract}

\section{Setup and the soft counterterm before the fix}

We consider the real emission of a gluon $i$ in a process with incoming partons
$1,2$, and set $s_{ab}=2\,k_a\cdot k_b$. Up to the overall normalisation and
colour conventions, which play no role here, the soft counterterm is a sum over
dipoles $\{l,m\}$ of eikonal factors times colour-correlated Born matrix
elements, each evaluated on a mapped Born configuration $\kb^{(i,l,m)}$:
\begin{equation}
  K_S^{(i)} \;=\; -\Nc \sum_{\{l,m\}} \E^{(i)}_{lm}\,
  B_{lm}\big(\kb^{(i,l,m)}\big),
  \qquad
  \E^{(i)}_{lm} = \frac{s_{lm}}{s_{il}\,s_{im}} .
  \label{eq:KS}
\end{equation}
For the dipole between the two incoming partons, the mapping $(i,a,b)$ with
emitter $a$ and spectator $b$ is the initial--initial one:
\begin{equation}
  x = \frac{s_{12}-s_{1i}-s_{2i}}{s_{12}},\qquad
  \kb_a = x\,k_a,\qquad \kb_b = k_b,\qquad
  \kb_j = \Lambda\,k_j \ \ (j\neq 1,2,i),
  \label{eq:FII}
\end{equation}
where $\Lambda$ is the Lorentz transformation taking
$K=k_1+k_2-k_i$ into $\bar K = \kb_1+\kb_2$. Both mappings give
$\bar K^2 = K^2 = x\,s_{12}$.

Decision D1 of the implementation was to average over the two beams:
\begin{equation}
  \E^{(i)}_{12}\,B_{12} \;\longrightarrow\;
  \frac12\,\E^{(i)}_{12}\Big[B_{12}\big(\kb^{(i,1,2)}\big)
                            + B_{12}\big(\kb^{(i,2,1)}\big)\Big].
  \label{eq:D1}
\end{equation}

\section{What goes wrong}

\paragraph{Collinear limit along beam 1.}
Take $k_i \to (1-z)\,k_1$ with $s_{1i}\to 0$ at fixed $z$. Then
$s_{2i}\to(1-z)\,s_{12}$, $x\to z$, and the eikonal is collinear singular:
\begin{equation}
  \E^{(i)}_{12} \;\to\; \frac{1}{(1-z)\,s_{1i}} .
\end{equation}
The two mapped Born configurations behave differently:
\begin{itemize}
\item mapping $(i,1,2)$: $\bar K \to z\,k_1 + k_2 = K$. This is the correct
      collinear Born, the same one on which the hard-collinear counterterm and
      its soft-collinear subtraction are evaluated;
\item mapping $(i,2,1)$: $\bar K \to k_1 + z\,k_2$. It has the same invariant
      mass $z\,s_{12}$, but it is related to the correct one by a longitudinal
      boost of rapidity
      \begin{equation}
        \Delta y = y\big(\kb^{(i,2,1)}\big) - y\big(\kb^{(i,1,2)}\big) = -\ln z .
      \end{equation}
\end{itemize}
The same holds for the other dipoles containing parton 1: the final--initial
and initial--final mappings give the correct collinear Born for both emitter
and spectator, so the issue is confined to the initial--initial dipole.

\paragraph{The leftover.}
Let $F$ be the observable, for instance the indicator function of a histogram
bin. In the limit $i\parallel 1$ the soft-collinear term removes
$\E_{12}B_{12}$ evaluated on the correct Born. With
eq.~\eqref{eq:D1}, what remains of the initial--initial dipole in the
subtracted integrand $R\,F(k) - K\,F(\kb)$ is
\begin{equation}
  \Delta \;\propto\; \frac{1}{2}\,\frac{1}{(1-z)\,s_{1i}}\;
  B_{12}\big(\kb^{(i,1,2)}\big)\,
  \Big[F\big(\kb^{(i,2,1)}\big) - F\big(\kb^{(i,1,2)}\big)\Big].
  \label{eq:leftover}
\end{equation}
The two Born configurations are related by a longitudinal boost, so all their
invariants coincide and the matrix element $B_{12}$ is the same in both
terms. The difference comes only from the observable.

\paragraph{Consequences.}
\begin{enumerate}
\item The IR-limit tests (\texttt{testR}) compare $R$ and $K$ without the
      observable. There the bracket in eq.~\eqref{eq:leftover} vanishes, so
      all limits pass.
\item For observables invariant under longitudinal boosts ($m_{ll}$, $p_T$,
      rapidity differences), $F$ is the same in both terms and $\Delta=0$.
\item Inclusively, $\sum_{\rm bins}F=1$ for both configurations: the bracket
      cancels in the sum over bins and total cross sections are correct.
\item Bin by bin in $y$ or $\eta$, the bracket is nonzero at fixed $z<1$, and
      the radiative integral contains $\int ds_{1i}/s_{1i}$. It diverges
      logarithmically. The $z$ integral is harmless, since the bracket vanishes
      like $\ln z\sim(1-z)$ as $z\to1$. In Monte Carlo the divergence shows up
      as an infinite-variance integrand: in the Drell--Yan run, per-bin errors
      of about $30$\,pb/unit on values of about $20$\,pb/unit in $y_{ll}$ and
      $\eta(e^-)$, while $m_{ll}$ and $p_T$ were clean.
\end{enumerate}
Neither the channel decomposition nor the damping factors help: the former
does not change where the counterterm pieces are filled, and $x^\alpha$ does not
vanish at finite $x$.

\paragraph{The single mapping of the paper.}
Using only $(i,2,1)$ for the initial--initial dipole, without D1, is correct
for $i\parallel2$. For $i\parallel1$, however, it puts the whole collinear
singularity on the boosted Born: the factor $\frac12$ in
eq.~\eqref{eq:leftover} becomes $1$. It is numerically worse
(table~\ref{tab:ppee}).

\section{The fix: partial fractions}

Decompose the eikonal into two fractions, each collinear singular along one
beam only:
\begin{equation}
  \E^{(i)}_{12} = \frac{s_{12}}{s_{1i}\,s_{2i}} = P_1 + P_2,
  \qquad
  P_a = \frac{s_{12}}{s_{ai}\,\big(s_{1i}+s_{2i}\big)} ,
  \label{eq:pf}
\end{equation}
and evaluate each on the mapping that rescales its own beam:
\begin{equation}
  \E^{(i)}_{12}\,B_{12} \;\longrightarrow\;
  P_1\,B_{12}\big(\kb^{(i,1,2)}\big) + P_2\,B_{12}\big(\kb^{(i,2,1)}\big).
  \label{eq:new}
\end{equation}

\paragraph{Limits.}
\begin{itemize}
\item $i\parallel1$: $P_1\to 1/[(1-z)s_{1i}]$ carries the full singularity on
      the correct Born, while $P_2\to 1/[(1-z)^2 s_{12}]$ is finite. The leftover is
      integrable bin by bin.
\item $i\parallel2$: the same, with $1\leftrightarrow2$.
\item $i$ soft: $x\to1$, both mapped configurations tend to the Born
      configuration, and $P_1+P_2=\E_{12}$. The soft limit is unchanged.
\end{itemize}
At a generic point, however, $K$ does change, because the two fractions
multiply different Born configurations. The singular limits are unchanged; the
values reached in the spurious limits (finite there) and the integrals of
single sectors change.

\paragraph{The integrated counterterm is unchanged.}
On the radiative phase space of mapping $(i,a,b)$, define $v$ by
\begin{equation}
  s_{ai} = (1-x)\,v\,s_{12},\qquad s_{bi}=(1-x)(1-v)\,s_{12},
\end{equation}
so that $\tilde v = s_{ai}/s_{12} = (1-x)\,v$ is the usual Catani--Seymour
variable. The $d$-dimensional measure is
\begin{equation}
  d\Phi_{\rm rad} \;\propto\; dx\;dv\;d\phi\,\sin^{-2\epsilon}\!\phi\;
  \big[(1-x)^2\,v(1-v)\big]^{-\epsilon}\,(1-x),
\end{equation}
which is symmetric under $v\leftrightarrow1-v$. In these variables
\begin{equation}
  P_a = \frac{1}{(1-x)^2\,s_{12}}\,\frac1v,
  \qquad
  \frac12\,\E_{12} = \frac{1}{(1-x)^2\,s_{12}}\,
  \frac{1}{2}\Big[\frac1v + \frac1{1-v}\Big].
\end{equation}
The mapped Born, the flux, the damping factor $x^\alpha$ and, after summing over
sectors with $\sum_j W_{ij}=1$, the whole integrand except the eikonal depend
on $x$ only. By the $v\leftrightarrow1-v$ symmetry of the measure,
\begin{equation}
  \int_0^1 dv\,\big[v(1-v)\big]^{-\epsilon}\,\frac1v
  \;=\;
  \int_0^1 dv\,\big[v(1-v)\big]^{-\epsilon}\,
  \frac12\Big[\frac1v+\frac1{1-v}\Big]
  \;=\; \frac{\Gamma(-\epsilon)\,\Gamma(1-\epsilon)}{\Gamma(1-2\epsilon)} ,
\end{equation}
so $\int P_a = \int\frac12\E_{12}$ exactly, for each beam. Each beam still
receives half of the integrated initial--initial soft dipole, which is what D1
was for, and the integrated counterterms $I^{\rm II}$ and $J^{\rm II}$
(2209.09123) are unchanged. Only the template of the local counterterm needs
to change.

\paragraph{Massive final states.}
The initial--initial dipole is between massless beams in $pp\to t\bar t$ as
well. The massive final-state dipoles and their mappings are not affected.

\section{Implementation}

In \texttt{tmp\_fortran/tmp\_files/NLO\_limits/S\_g.f}, the loop over the
dipoles $m<l$ uses the mapping $(i,k_e,k_r)$ with emitter $k_e$ and recoiler
$k_r$. For the initial--initial dipole ($l\le2$) both mappings are used, with
weight $w = s_{i k_r}/(s_{il}+s_{im})$, since
$\E_{lm}\,w = P_{k_e}$:
\begin{lstlisting}
            nmap=1
            if(l.le.2)nmap=2
            wmap=1d0
            do imap=1,nmap
            if(imap.eq.1)then
               ke=l
               kr=m
            else
               ke=m
               kr=l
            endif
            if(nmap.eq.2)wmap=xs(i,kr)/(xs(i,l)+xs(i,m))
            ...
            M2tmp=M2tmp*damp*xj*wmap
\end{lstlisting}
Before the fix the weight was \texttt{wmap=1d0/dble(nmap)}, i.e.\ $\frac12$.

\section{Validation}

All runs are in the cloud copy, with processes regenerated from the new
template, default statistics and the same random seed.

\paragraph{$pp\to e^+e^-$.}
\texttt{testR} passes in all sectors. Table~\ref{tab:ppee} compares the median
per-bin Monte Carlo error of the real-emission contribution.

\begin{table}[h]
\centering
\begin{tabular}{lccc}
\toprule
 & D1 (average) & single $(i,2,1)$ & partial fractions \\
\midrule
$\sigma_R$ [pb]                    & 128.3644 & 128.3457 & 128.3656 \\
median error, $y_{ll}$ bins        & 33.9     & 54.7     & 1.52 \\
median error, $\eta(e^-)$ bins     & 36.1     & 45.2     & 4.07 \\
$m_{ll}$, $p_{T,ll}$               & \multicolumn{3}{c}{unchanged} \\
\bottomrule
\end{tabular}
\caption{$pp\to e^+e^-$, real-emission contribution: total and median per-bin
error of the distributions, for the three choices of the initial--initial
soft dipole.}
\label{tab:ppee}
\end{table}

\paragraph{$pp\to Zj$.}
\texttt{testR}: 79 sectors OK and 3 flagged, none of them a problem.
\begin{itemize}
\item $gd\to Zgd$ and $gu\to Zgu$, sector 4\_5, spurious limit $C_{jr}$
      flagged as ``grows $\times84$'' (old template: $\times3.0$). On all 20
      test events $K$ converges: the relative change in the last step is below
      $10^{-4}$, and $|K|$ at small $\lambda$ is at most $1.03$ times its
      value at large $\lambda$. It converges to a different finite value
      (worst event: $0.310$ instead of $0.209$). The flag is an artefact of the
      growth measure, which divides by $K$ at $\lambda=10^{-2}$, where $K$
      happens to pass close to zero.
\item $d\bar u\to Zd\bar u$, sector 4\_1, $C_{ij}$ at $1.1\times10^{-3}$
      against a tolerance of $10^{-3}$. This process has no gluon, so the
      sector contains no soft-gluon code and is unaffected by the change.
\end{itemize}
Real-emission integral: $-9242\pm167$\,pb with partial fractions, and
$-9293\pm176$\,pb with the old weight restored in the 24 sectors that contain
an initial--initial soft gluon. The other 58 sectors are bit-identical. Per
process the totals agree (pulls $-0.50$, $-1.02$, $+0.73$, $+1.74$, $+0.54$,
$-0.91$, so $\chi^2=5.97$ for 6 processes), with errors of the same size. Single sectors may
differ, since only the sum over sectors carries the invariant soft integral.

\paragraph{$pp\to t\bar t$.}
The regenerated process differs from the previous one only in the soft block
of the 8 sectors with an initial--initial dipole. \texttt{testR} passes 12/12.
$\sigma_R = -79.510\pm0.288$\,pb with partial fractions, against
$-79.516\pm0.290$\,pb before. The 4 unaffected sectors are bit-identical and
the per-process pulls are at most $0.02$.

\end{document}
